\documentclass[12pt]{article}

% Idioma y codificación
\usepackage[spanish, es-tabla, es-notilde]{babel}       %es-tabla para que se titule "Tabla"
\usepackage[utf8]{inputenc}

% Márgenes
\usepackage[a4paper,top=3cm,bottom=2.5cm,left=3cm,right=3cm]{geometry}

% Comentarios de bloque
\usepackage{verbatim}

% Paquetes de links
\usepackage[hidelinks]{hyperref}    % Permite enlaces
\usepackage{url}                    % redirecciona a la web

% Más opciones para enumeraciones
\usepackage{enumitem}

% Personalizar la portada
\usepackage{titling}

% Paquetes de tablas
\usepackage{multirow}

% Para añadir el símbolo de euro
\usepackage{eurosym}

% Paquete para subrayar con diferentes estilos
\usepackage[normalem]{ulem}


%------------------------------------------------------------------------

%Paquetes de figuras
\usepackage{caption}
\usepackage{subcaption} % Figuras al lado de otras
\usepackage{float}      % Poner figuras en el sitio indicado H.


% Paquetes de imágenes
\usepackage{graphicx}       % Paquete para añadir imágenes
\usepackage{transparent}    % Para manejar la opacidad de las figuras

% Paquete para usar colores
\usepackage[dvipsnames, table, xcdraw]{xcolor}
\usepackage{pagecolor}      % Para cambiar el color de la página

% Habilita tamaños de fuente mayores
\usepackage{fix-cm}

% Para los gráficos
\usepackage{tikz}
\usepackage{forest}

% Para poder situar los nodos en los grafos
\usetikzlibrary{positioning}

% Para evitar problemas entre babel y tikz
\usetikzlibrary{babel}


%------------------------------------------------------------------------

% Paquetes de matemáticas
\usepackage{mathtools, amsfonts, amssymb, mathrsfs}
\usepackage[makeroom]{cancel}     % Simplificar tachando
\usepackage{polynom}    % Divisiones y Ruffini
\usepackage{units} % Para poner fracciones diagonales con \nicefrac

\usepackage{pgfplots}   %Representar funciones
\pgfplotsset{compat=1.18}  % Versión 1.18

\usepackage{tikz-cd}    % Para usar diagramas de composiciones
\usetikzlibrary{calc}   % Para usar cálculo de coordenadas en tikz

%Definición de teoremas, etc.
\usepackage{amsthm}
%\swapnumbers   % Intercambia la posición del texto y de la numeración

\theoremstyle{plain}

\makeatletter
\@ifclassloaded{article}{
  % --- Configuración para ARTICLE ---
  \newtheorem{teo}{Teorema}[section]
  \theoremstyle{definition}
  \newtheorem{definicion}{Definición}[section]
  \newtheorem{ejercicio}{Ejercicio}[section]

  \renewcommand{\theteo}{%
    \ifnum\value{section}=0 \arabic{teo}\else \thesection.\arabic{teo}\fi
  }
  \renewcommand{\thedefinicion}{%
    \ifnum\value{section}=0 \arabic{definicion}\else \thesection.\arabic{definicion}\fi
  }
  \renewcommand{\theejercicio}{%
    \ifnum\value{section}=0 \arabic{ejercicio}\else \thesection.\arabic{ejercicio}\fi
  }
}{
  % --- Configuración para BOOK / REPORT ---
  \newtheorem{teo}{Teorema}[chapter]
  \theoremstyle{definition}
  \newtheorem{definicion}{Definición}[chapter]
  \newtheorem{ejercicio}{Ejercicio}[section]

  % Si no hay sección: Capítulo.Teorema (ej. 2.1)
  % Si hay sección:    Capítulo.Sección.Teorema (ej. 2.3.1)
  \renewcommand{\theteo}{%
    \ifnum\value{section}=0
      \thechapter.\arabic{teo}%
    \else
      \thesection.\arabic{teo}% (\thesection ya contiene el número de capítulo)
    \fi
  }
  \renewcommand{\thedefinicion}{%
    \ifnum\value{section}=0
      \thechapter.\arabic{definicion}%
    \else
      \thesection.\arabic{definicion}%
    \fi
  }
  \renewcommand{\theejercicio}{%
    \ifnum\value{section}=0
      \arabic{ejercicio}%
    \else
      \thesection.\arabic{ejercicio}% Formato por defecto: ej. 2.2.1
    \fi
  }
}
\makeatother

% Teoremas dependientes o con contador compartido
\newtheorem{coro}{Corolario}[teo]
\newtheorem{prop}[teo]{Proposición}
\newtheorem{lema}[teo]{Lema}
\newtheorem*{lema*}{Lema}

\theoremstyle{remark}
\newtheorem*{observacion}{Observación}

\theoremstyle{definition}
\newtheorem*{notacion}{Notación}
\newtheorem*{ejemplo}{Ejemplo}
\newtheorem*{ejercicio*}{Ejercicio}


% \renewcommand\qedsymbol{$\blacksquare$}         % Cambiar símbolo QED
%------------------------------------------------------------------------

% Paquetes para encabezados
\usepackage{fancyhdr}
\pagestyle{fancy}
\fancyhf{}

\newcommand{\helv}{ % Modificación tamaño de letra
\fontfamily{}\fontsize{12}{12}\selectfont}
\setlength{\headheight}{15pt} % Amplía el tamaño del índice


%\usepackage{lastpage}   % Referenciar última pag   \pageref{LastPage}
%\fancyfoot[C]{%
%  \begin{minipage}{\textwidth}
%    \centering
%    ~\\
%    \thepage\\
%    \href{https://losdeldgiim.github.io/}{\texttt{\footnotesize losdeldgiim.github.io}}
%  \end{minipage}
%}
\fancyfoot[C]{\thepage}
\fancyfoot[R]{\href{https://losdeldgiim.github.io/}{\texttt{\footnotesize losdeldgiim.github.io}}}

%------------------------------------------------------------------------

% Conseguir que no ponga "Capítulo 1". Sino solo "1."
\makeatletter
\@ifclassloaded{book}{
  \renewcommand{\chaptermark}[1]{\markboth{\thechapter.\ #1}{}} % En el encabezado
    
  \renewcommand{\@makechapterhead}[1]{%
  \vspace*{50\p@}%
  {\parindent \z@ \raggedright \normalfont
    \ifnum \c@secnumdepth >\m@ne
      \huge\bfseries \thechapter.\hspace{1em}\ignorespaces
    \fi
    \interlinepenalty\@M
    \Huge \bfseries #1\par\nobreak
    \vskip 40\p@
  }}
}
\makeatother

%------------------------------------------------------------------------
% Paquetes de cógido
\usepackage{minted}
\renewcommand\listingscaption{Código fuente}

\usepackage{fancyvrb}
% Personaliza el tamaño de los números de línea
\renewcommand{\theFancyVerbLine}{\small\arabic{FancyVerbLine}}

% Estilo para C++
\newminted{cpp}{
    frame=lines,
    framesep=2mm,
    baselinestretch=1.2,
    linenos,
    escapeinside=||
}

% para minted
\definecolor{LightGray}{rgb}{0.95,0.95,0.92}
\setminted{
    linenos=true,
    stepnumber=5,
    numberfirstline=true,
    autogobble,
    breaklines=true,
    breakautoindent=true,
    breaksymbolleft=,
    breaksymbolright=,
    breaksymbolindentleft=0pt,
    breaksymbolindentright=0pt,
    breaksymbolsepleft=0pt,
    breaksymbolsepright=0pt,
    fontsize=\footnotesize,
    bgcolor=LightGray,
    numbersep=10pt
}


\usepackage{listings} % Para incluir código desde un archivo

\renewcommand\lstlistingname{Código Fuente}
\renewcommand\lstlistlistingname{Índice de Códigos Fuente}

% Definir colores
\definecolor{vscodepurple}{rgb}{0.5,0,0.5}
\definecolor{vscodeblue}{rgb}{0,0,0.8}
\definecolor{vscodegreen}{rgb}{0,0.5,0}
\definecolor{vscodegray}{rgb}{0.5,0.5,0.5}
\definecolor{vscodebackground}{rgb}{0.97,0.97,0.97}
\definecolor{vscodelightgray}{rgb}{0.9,0.9,0.9}

% Configuración para el estilo de C similar a VSCode
\lstdefinestyle{vscode_C}{
  backgroundcolor=\color{vscodebackground},
  commentstyle=\color{vscodegreen},
  keywordstyle=\color{vscodeblue},
  numberstyle=\tiny\color{vscodegray},
  stringstyle=\color{vscodepurple},
  basicstyle=\scriptsize\ttfamily,
  breakatwhitespace=false,
  breaklines=true,
  captionpos=b,
  keepspaces=true,
  numbers=left,
  numbersep=5pt,
  showspaces=false,
  showstringspaces=false,
  showtabs=false,
  tabsize=2,
  frame=tb,
  framerule=0pt,
  aboveskip=10pt,
  belowskip=10pt,
  xleftmargin=10pt,
  xrightmargin=10pt,
  framexleftmargin=10pt,
  framexrightmargin=10pt,
  framesep=0pt,
  rulecolor=\color{vscodelightgray},
  backgroundcolor=\color{vscodebackground},
}

%------------------------------------------------------------------------
% La etiqueda IA sirve para indicar a partir de donde le pasamos el preambulo a la IA para el escaneo
% IA

% Comandos definidos
\newcommand{\bb}[1]{\mathbb{#1}}
\newcommand{\cc}[1]{\mathcal{#1}}

% I prefer the slanted \leq
\let\oldleq\leq % save them in case they're every wanted
\let\oldgeq\geq
\renewcommand{\leq}{\leqslant}
\renewcommand{\geq}{\geqslant}

% Si y solo si
\newcommand{\sii}{\iff}

% MCD y MCM
\DeclareMathOperator{\mcd}{mcd}
\DeclareMathOperator{\mcm}{mcm}

% Signo
\DeclareMathOperator{\sgn}{sgn}

% Imagen
\DeclareMathOperator{\Ima}{Im}

% Aplicación identidad
\newcommand{\id}{\mathrm{Id}}

% Letras griegas
\newcommand{\eps}{\epsilon}
\newcommand{\veps}{\varepsilon}
\newcommand{\lm}{\lambda}

\newcommand{\ol}{\overline}
\newcommand{\ul}{\underline}
\newcommand{\wt}{\widetilde}
\newcommand{\wh}{\widehat}

\let\oldvec\vec
\renewcommand{\vec}{\overrightarrow}

% Derivadas parciales
\newcommand{\del}[2]{\frac{\partial #1}{\partial #2}}
\newcommand{\Del}[3]{\frac{\partial^{#1} #2}{\partial #3^{#1}}}
\newcommand{\deld}[2]{\dfrac{\partial #1}{\partial #2}}
\newcommand{\Deld}[3]{\dfrac{\partial^{#1} #2}{\partial #3^{#1}}}


\newcommand{\AstIg}{\stackrel{(\ast)}{=}}
\newcommand{\Hop}{\stackrel{L'H\hat{o}pital}{=}}

\newcommand{\red}[1]{{\color{red}#1}} % Para integrales, destacar los cambios.

% Método de integración
\newcommand{\MetInt}[2]{
    \left[\begin{array}{c}
        #1 \\ #2
    \end{array}\right]
}

% Declarar aplicaciones
% 1. Nombre aplicación
% 2. Dominio
% 3. Codominio
% 4. Variable
% 5. Imagen de la variable
\newcommand{\Func}[5]{
    \begin{equation*}
        \begin{array}{rrll}
            \displaystyle #1:& \displaystyle  #2 & \longrightarrow & \displaystyle  #3\\
               & \displaystyle  #4 & \longmapsto & \displaystyle  #5
        \end{array}
    \end{equation*}
}

%------------------------------------------------------------------------


\newcommand{\myhref}[2]{\href{#1}{\color{blue}\uline{#2}}}
\newcommand{\T}[0]{\cc{T}}

\begin{document}

    % 1. Foto de fondo
    % 2. Título
    % 3. Encabezado Izquierdo
    % 4. Color de fondo
    % 5. Coord x del titulo
    % 6. Coord y del titulo
    % 7. Fecha

    
    % 1. Foto de fondo
% 2. Título
% 3. Encabezado Izquierdo
% 4. Color de fondo
% 5. Coord x del titulo
% 6. Coord y del titulo
% 7. Fecha
% 8. Autor

\newcommand{\portada}[8]{

    \portadaBase{#1}{#2}{#3}{#4}{#5}{#6}{#7}{#8}
    \portadaBook{#1}{#2}{#3}{#4}{#5}{#6}{#7}{#8}
}

\newcommand{\portadaExamen}[8]{

    \portadaBase{#1}{#2}{#3}{#4}{#5}{#6}{#7}{#8}
    \portadaArticle{#1}{#2}{#3}{#4}{#5}{#6}{#7}{#8}
}




\newcommand{\portadaBase}[8]{

    % Tiene la portada principal y la licencia Creative Commons
    
    % 1. Foto de fondo
    % 2. Título
    % 3. Encabezado Izquierdo
    % 4. Color de fondo
    % 5. Coord x del titulo
    % 6. Coord y del titulo
    % 7. Fecha
    % 8. Autor
    
    
    \thispagestyle{empty}               % Sin encabezado ni pie de página
    \newgeometry{margin=0cm}        % Márgenes nulos para la primera página
    
    
    % Encabezado
    \fancyhead[L]{\helv #3}
    \fancyhead[R]{\helv \nouppercase{\leftmark}}
    
    
    \pagecolor{#4}        % Color de fondo para la portada
    
    \begin{figure}[p]
        \centering
        \transparent{0.3}           % Opacidad del 30% para la imagen
        
        \includegraphics[width=\paperwidth, keepaspectratio]{../../_assets/#1}
    
        \begin{tikzpicture}[remember picture, overlay]
            \node[anchor=north west, text=white, opacity=1, font=\fontsize{60}{90}\selectfont\bfseries\sffamily, align=left] at (#5, #6) {#2};
            
            \node[anchor=south east, text=white, opacity=1, font=\fontsize{12}{18}\selectfont\sffamily, align=right] at (9.7, 3) {\href{https://losdeldgiim.github.io/}{\textbf{Los Del DGIIM}, \texttt{\footnotesize losdeldgiim.github.io}}};
            
            \node[anchor=south east, text=white, opacity=1, font=\fontsize{12}{15}\selectfont\sffamily, align=right] at (9.7, 1.8) {Doble Grado en Ingeniería Informática y Matemáticas\\Universidad de Granada};
        \end{tikzpicture}
    \end{figure}
    
    
    \restoregeometry        % Restaurar márgenes normales para las páginas subsiguientes
    \pagecolor{white}       % Restaurar el color de página
    
    
    \newpage
    \thispagestyle{empty}               % Sin encabezado ni pie de página
    \begin{tikzpicture}[remember picture, overlay]
        \node[anchor=south west, inner sep=3cm] at (current page.south west) {
            \begin{minipage}{0.5\paperwidth}
                \href{https://creativecommons.org/licenses/by-nc-nd/4.0/}{
                    \includegraphics[height=2cm]{../../_assets/Licencia.png}
                }\vspace{1cm}\\
                Esta obra está bajo una
                \href{https://creativecommons.org/licenses/by-nc-nd/4.0/}{
                    Licencia Creative Commons Atribución-NoComercial-SinDerivadas 4.0 Internacional (CC BY-NC-ND 4.0).
                }\\
    
                Eres libre de compartir y redistribuir el contenido de esta obra en cualquier medio o formato, siempre y cuando des el crédito adecuado a los autores originales y no persigas fines comerciales. 
            \end{minipage}
        };
    \end{tikzpicture}
    
    
    
    % 1. Foto de fondo
    % 2. Título
    % 3. Encabezado Izquierdo
    % 4. Color de fondo
    % 5. Coord x del titulo
    % 6. Coord y del titulo
    % 7. Fecha
    % 8. Autor


}


\newcommand{\portadaBook}[8]{

    % 1. Foto de fondo
    % 2. Título
    % 3. Encabezado Izquierdo
    % 4. Color de fondo
    % 5. Coord x del titulo
    % 6. Coord y del titulo
    % 7. Fecha
    % 8. Autor

    % Personaliza el formato del título
    \pretitle{\begin{center}\bfseries\fontsize{42}{56}\selectfont}
    \posttitle{\par\end{center}\vspace{2em}}
    
    % Personaliza el formato del autor
    \preauthor{\begin{center}\Large}
    \postauthor{\par\end{center}\vfill}
    
    % Personaliza el formato de la fecha
    \predate{\begin{center}\huge}
    \postdate{\par\end{center}\vspace{2em}}
    
    \title{#2}
    \author{\href{https://losdeldgiim.github.io/}{Los Del DGIIM, \texttt{\large losdeldgiim.github.io}}
    \\ \vspace{0.5cm}#8}
    \date{Granada, #7}
    \maketitle
    
    \tableofcontents
}




\newcommand{\portadaArticle}[8]{

    % 1. Foto de fondo
    % 2. Título
    % 3. Encabezado Izquierdo
    % 4. Color de fondo
    % 5. Coord x del titulo
    % 6. Coord y del titulo
    % 7. Fecha
    % 8. Autor

    % Personaliza el formato del título
    \pretitle{\begin{center}\bfseries\fontsize{42}{56}\selectfont}
    \posttitle{\par\end{center}\vspace{2em}}
    
    % Personaliza el formato del autor
    \preauthor{\begin{center}\Large}
    \postauthor{\par\end{center}\vspace{3em}}
    
    % Personaliza el formato de la fecha
    \predate{\begin{center}\huge}
    \postdate{\par\end{center}\vspace{5em}}
    
    \title{#2}
    \author{\href{https://losdeldgiim.github.io/}{Los Del DGIIM, \texttt{\large losdeldgiim.github.io}}
    \\ \vspace{0.5cm}#8}
    \date{Granada, #7}
    \thispagestyle{empty}               % Sin encabezado ni pie de página
    \maketitle
    \vfill
}
    \portadaExamen{ffccA4.jpg}{Topología I\\Examen XVI}{Topología I. Examen XVI}{MidnightBlue}{-8}{28}{2026}{Lucía Benjumeda Maldonado}

    \begin{description}
        \item[Asignatura] Topología I.
        \item[Curso Académico] 2026-27.
        \item[Grado] Grado en Matemáticas.
        \item[Grupo] B.
        \item[Profesor] Manuel César Rosales Lombardo.
        \item[Descripción] Parcial.
        \item[Fecha] 19 de noviembre de 2026.
        % \item[Duración] 60 minutos.
    
    \end{description}
    \newpage

    \begin{ejercicio}
       En $\bb{R}$, se considera la topología $T$ de la que conocemos una base dada por:
        \[
        \mathcal{B}=\{[-r,r]\quad : \quad r>0\}.
        \]
        \begin{enumerate}
            \item Estudiar si $(\mathbb{R}, T)$ es un espacio $T_1$. ¿Es cierto que $T\cap C_T = \{\emptyset, \mathbb{R}\}$?\\
            Veamos que $(\mathbb{R}, T)$ no es un espacio $T_1$. Se puede ver claramente tomando $x = 0$ e $y \in \mathbb{R} \setminus \{0\}$ (por ejemplo, $y \in \mathbb{R}^+$).
            Para cualquier abierto $U_y \in T$ tal que $y \in U_y$, como $\mathcal{B}$ es base de $T$, existe algún básico $[-r, r] \in \mathcal{B}$ (con $r > 0$) tal que $y \in [-r, r] \subseteq U_y$. 
            Sin embargo, como $r > 0$, se cumple que $0 \in [-r, r] \subseteq U_y$. Es decir, $0$ está en todo abierto no vacío de $T$. Por tanto, todo punto distinto de $0$ comparte todos sus entornos con el $0$, luego $(\mathbb{R}, T)$ no es $T_1$ (de hecho, ni siquiera es $T_0$).\\

            Para probar que $T\cap C_T = \{\emptyset, \mathbb{R}\}$, lo haremos por doble inclusión:
            \begin{itemize}
                \item[$\supseteq$] Es evidente que $\emptyset, \mathbb{R} \in T \cap C_T$ por definición de topología.
                \item[$\subseteq$] Ahora veamos que sea $U \in T \cap C_T$ con $U \neq \emptyset$, probemos que $U = \mathbb{R}$.\\
                Como $U \in T$ y $U \neq \emptyset$, es unión de elementos de $\mathcal{B}$ y, como hemos visto antes, todo básico contiene al $0$, luego $0 \in U$.\\
                Por otro lado, como $U \in C_T$, su complementario $V = \mathbb{R} \setminus U$ es abierto. Si $V \neq \emptyset$, al ser también unión de elementos de $\mathcal{B}$, tendríamos que $0 \in V = \mathbb{R} \setminus U$ lo cual es una contradicción porque $0\in U \Rightarrow V = \emptyset$, de donde concluimos que $U = \mathbb{R}$. 
            \end{itemize}
        \item Para cada $x\neq 0$, encontrar una base de entornos de $x$ en $(\mathbb{R}, T)$ con un único elemento.
        La base de entornos más económica (de cardinal 1) viene dada por:
        \[
            \mathcal{B}_x = \{[-x, x]\} \text{ si } x > 0, \qquad \mathcal{B}_x = \{[x, -x]\} \text{ si } x < 0
        \]
        (es decir, $\mathcal{B}_x = \{[-|x|, |x|]\}$ para todo $x \in \bb{R} \setminus \{0\}$).
        \\Probemos que $\mathcal{B}_x$ es efectivamente una base de entornos para cada $x \in \bb{R} \setminus \{0\}$:
        \begin{itemize}
            \item \textbf{Reflexividad puntual:} $\mathcal{B}_x \neq \emptyset$ y para el único elemento $V = [-|x|, |x|] \in \mathcal{B}_x$ se cumple claramente que $x \in [-|x|, |x|] = V$.
            
            \item \textbf{Intersección finita:} Sean $V_1, V_2 \in \mathcal{B}_x$. Como $\mathcal{B}_x$ es unitaria, necesariamente $V_1 = V_2 = [-|x|, |x|]$. Tomando $V_3 = [-|x|, |x|] \in \mathcal{B}_x$, se verifica trivialmente que $V_3 \subseteq V_1 \cap V_2$.
            
            \item \textbf{Axioma de Bourbaki:} Sea $V = [-|x|, |x|] \in \mathcal{B}_x$. Como $|x| > 0$, tenemos que $V \in \mathcal{B} \subseteq T$ es decir, $V$ ya es abierto, luego entorno de todos sus puntos por un entorno básico ($\forall y\in V \quad \exists V_y\in \mathcal{B}_y$ tal que  $y\in V_y\in V$)
        \end{itemize}
        Probado así que $\mathcal{B}_x$ conforma una base de entornos para cada $x \neq 0$.
                
        \item Demostrar que $(\mathbb{R}, T)$ cumple el primer axioma de numerabilidad pero no el segundo.
        \\Teniendo definida la base de entornos $\forall x\in \mathbb{R}\setminus\{0\}$, agregemos también para $x=0$, de manera que:
        \[\begin{cases}
            \mathcal{B}_x = \{[-|x|, |x|]\} \quad\text{para todo } x \in \bb{R} \setminus \{0\}\\
            \mathcal{B}_0 = \{0\} \quad\text{que es en realidad un caso particular donde }\{0\}=[-0,0]
        \end{cases}\]
        Luego $\mathcal{B}_x$ es finita $\forall x\in \bb{R}$ y, por tanto, numerable.\\

        Ahora veremos que no cumple el segundo axioma de numerabilidad. Supongamos por reducción al absurdo que $(\bb{R}, T)$ es (2AN). Como $\mathcal{B}$ es una base de la topología $T$, por ser (2AN) existirá una subfamilia numerable $\mathcal{B}_n \subseteq \mathcal{B}$ que también es base de $T$:
        \[
            \mathcal{B}_n = \{[-r_n, r_n] : n \in \bb{N}\} \quad \text{con } r_n \in \bb{R}^+.
        \]
        Sea $R = \{r_n : n \in \bb{N}\} \subseteq \bb{R}^+$. Como $R$ es numerable y $\bb{R}^+$ es no numerable, se tiene que $\bb{R}^+ \setminus R \neq \emptyset$.
        
        Tomemos $x \in \bb{R}^+ \setminus R \Rightarrow [-x, x] \in \mathcal{B} \subseteq T$, por lo que es abierto y debe expresarse como unión de elementos de $\mathcal{B}_N$:
        \[
            [-x, x] = \bigcup_{n \in J} [-r_n, r_n] \quad \text{para algún } J \subseteq \bb{N}.
        \]
        Como $x \in [-x, x]$, existe algún $n_0 \in J$ tal que $x \in [-r_{n_0}, r_{n_0}]$, lo que implica que $x \le r_{n_0}$. Por otro lado, como $[-r_{n_0}, r_{n_0}] \subseteq [-x, x]$, también se cumple que $r_{n_0} \le x$. De ambas desigualdades deducimos que $r_{n_0} = x$, lo cual es una contradicción pues $x \notin R$. Por tanto, no existe base numerable y $(\bb{R}, T)$ NO es (2AN).
            
        \item Calcular el cierre en $(\mathbb{R}, T)$ del conjunto $A=\{1\}$ y el interior del conjunto $C=(-3,5)$
        \begin{itemize}
    \item \textbf{Probemos que $\overline{A} = \bb{R} \setminus ]-1, 1[$:}
    \begin{itemize}
        \item[$\subseteq$)] Probemos por contrarrecíproco que si $x \in ]-1, 1[ = \mathbb{R}\setminus(\bb{R}\setminus ]-1,1[)$, entonces $x \notin \overline{A}$ (es decir, que ningún punto fuera de $\bb{R} \setminus ]-1, 1[$ pertenece a la adherencia)\\
        Si $x \in ]-1, 1[ \setminus \{0\}$, sabemos que $|x| < 1$. Considerando su entorno básico $[-|x|, |x|] \in \mathcal{B}_x$, es evidente que $[-|x|, |x|] \cap \{1\} = \emptyset$ pues $|x| < 1$. Por tanto, $x \notin \overline{A}$. \\
        Por contrarrecíproco, todo $x \in \overline{A}$ cumple que $x \in \bb{R} \setminus ]-1, 1[$.

        \item[$\supseteq$)] Sea $x \in \bb{R} \setminus ]-1, 1[$ arbitrario, es decir, $|x| \ge 1$. Como $x \neq 0$, su base de entornos es $\mathcal{B}_x = \{[-|x|, |x|]\}$. Dado que $|x| \ge 1$, se tiene que $[-1, 1] \subseteq [-|x|, |x|]$, por lo que:
        \[
            [-|x|, |x|] \cap \{1\} = \{1\} \neq \emptyset.
        \]
        Como el único entorno básico de $x$ corta a $A$, todo entorno de $x$ corta a $A$, de donde $x \in \overline{A}$.
    \end{itemize}

    \item \textbf{Ahora veamos que $C^\circ = ]-3, 3[$:}
    \begin{itemize}
        \item[$\subseteq$)] Probemos por contrarrecíproco que si $x \in \bb{R} \setminus ]-3, 3[$, entonces $x \notin C^\circ$. \\
        Sea $x \in \bb{R} \setminus ]-3, 3[$, es decir, $|x| \ge 3$. Como $x \neq 0$, su único entorno básico es $[-|x|, |x|]$. Dado que $|x| \ge 3 \Rightarrow [-|x|, |x|]\nsubseteq C$, lo que equivale a probar que $[-|x|, |x|]\cap(\bb{R}\setminus C)\neq \emptyset$.\\
        Efectivamente $-3\in [-|x|, |x|]\cap(\bb{R}\setminus C) \Rightarrow [-|x|, |x|] \nsubseteq C$, lo que implica que ningún entorno de $x$ está contenido en $C$, luego $x \notin C^\circ$. Por contrarrecíproco, $C^\circ \subseteq ]-3, 3[$.

        \item[$\supseteq$)] Sea $x \in ]-3, 3[$. Si $x = 0$, el básico $[-1, 1]$ es un entorno de $0$ contenido en $(-3, 5)$, luego $0 \in C^\circ$. \\
        Si $x \in ]-3, 3[ \setminus \{0\}$, entonces $|x| < 3$, por lo que su entorno básico cumple:
        \[
            x \in [-|x|, |x|] \subseteq ]-3, 3[ \subseteq (-3, 5) = C,
        \]
        de donde concluimos directamente que $x \in C^\circ$.
    \end{itemize}
\end{itemize}
        \item Probar que si un subconjunto $D\subseteq \mathbb{R}$ es discreto, entonces consta de un único punto.\\
        Veamos que si $(D, T_{|D}) = (D, T_d)$ , por definición de topología inducida tendremos que:
        \[
            T_{|D} = \{U \cap D : U \in T\} = \left\{ U \cap D : U = \bigcup_{i \in I} B_i \text{ con } B_i \in \mathcal{B} \right\}
        \]
        Y por tanto su base es $\mathcal{B}_D = \{B \cap D : B \in \mathcal{B}\} = \{[-r, r] \cap D : r > 0\}$.
        
        Como $(D, T_{|D})$ es discreto, todo subconjunto unitario $\{x\}\in D$ es abierto en $T_{|D}$, por lo que debe expresarse como unión de elementos de $\mathcal{B}_D$ (en particular, existe algún $r > 0$ tal que $\{x\} = [-r, r] \cap D$).
        
        Supongamos por reducción al absurdo que $D$ contiene al menos dos puntos distintos $x, y \in D$ ($x \neq y$). Como $\{x\} \in T_{|D}$, existe un básico $[-r_1, r_1] \in \mathcal{B}$ (con $r_1 > 0$) tal que:
        \[
            \{x\} = [-r_1, r_1] \cap D.
        \]
        Como $x \in [-r_1, r_1]$, se tiene que $|x| \le r_1$. Además, como $y \in D$ e $y \neq x$, necesariamente $y \notin [-r_1, r_1]$, por lo que $|y| > r_1$. Combinando ambas desigualdades obtenemos:
        \[
            |x| \le r_1 < |y| \implies |x| < |y|.
        \]
        Ahora, como $\{y\} \in T_{|D}$, existe $r_2 > 0$ tal que $\{y\} = [-r_2, r_2] \cap D$), por un razonamiento totalmente análogo se deduce que:
        \[
            |y| \le r_2 < |x| \implies |y| < |x|,
        \]
        llegando a la contradicción $|x| < |y| < |x|$. Por tanto, $D$ no puede tener más de un punto.
        \end{enumerate}
    \end{ejercicio}
\end{document}